\begin{helper}[Product Fiber Sum]
Fix $n\in\mathbb N$, $I\subseteq\mathrm{Fin}(n)$, parameter $M\in\mathbb N$, and $p_0\in\mathbb R$ with $0\le p_0\le1$.
Let $\pi:\mathrm{Fin}(n)\hookrightarrow\mathrm{Fin}(M)\times\mathrm{Fin}(M)$ be an embedding.
Let $P_R=\{\pi(u)_1:u\in I\}$, $P_B=\{\pi(u)_2:u\in I\}$.
Let $q=\noderef{FixedSetProductModelMassFn}(p_0,P_R,P_B)$.
Let $\mathcal D$ be the finite set of incident-data families $D=(D_R, D_B)$ where $D_R:\mathrm{Fin}(M)\setminus P_R\to\mathcal{P}(P_R)$ and $D_B:\mathrm{Fin}(M)\setminus P_B\to\mathcal{P}(P_B)$.
For a product assignment $v$, define its incident data $D(v)$ by $D_R(x)=v_x^{(1)}$ for $x\notin P_R$, and $D_B(y)=v_{y+M}^{(2)}$ for $y\notin P_B$.
Then for any predicate $H$ on the data, we have
\[
\sum_{\{v:H(D(v))\}}\prod_iq(v_i) = \sum_{\{D\in\mathcal D: H(D)\}} \nu(D),
\]
where $\nu(D)$ is the mass
\[
\prod_{x\notin P_R} p_0^{|D_R(x)|}(1-p_0)^{|P_R|-|D_R(x)|} \prod_{y\notin P_B} p_0^{|D_B(y)|}(1-p_0)^{|P_B|-|D_B(y)|}.
\]
\end{helper}
\begin{proof}
For data \(D\), let \(F_D=\{v:D(v)=D\}\) be the finite fiber of product
assignments.  The fibers \(F_D\) partition the full product assignment space,
so finite distributivity gives
\[
\sum_{\{v:H(D(v))\}}\prod_i q(v_i)
=
\sum_D {\bf 1}_{H(D)}
   \sum_{v\in F_D}\prod_iq(v_i).
\]
It remains to compute the inner fiber mass.

First consider invalid data.  If \(D_R(x)\ne\emptyset\) for some \(x\in P_R\)
or \(D_B(y)\ne\emptyset\) for some \(y\in P_B\), then \(F_D=\emptyset\),
because \(D(v)\) is defined to be \(\emptyset\) on those embedded
coordinates.  If \(x\notin P_R\) but \(D_R(x)\not\subseteq P_R\), then every
assignment in the fiber has first component \(D_R(x)\) at that red coordinate,
and \(\noderef{FixedSetProductModelMassFn}\) assigns that coordinate weight
\(0\).  The same applies if \(D_B(y)\not\subseteq P_B\).  Hence invalid data
have fiber mass \(0\).

Assume now that \(D\) is valid.  Expand the product over the \(2M\) product
coordinates and apply finite Fubini, so the fiber mass factors into one
coordinate sum for each coordinate.  For a red coordinate \(x\notin P_R\), the
fiber fixes the first component to \(R=D_R(x)\) and leaves the second
component arbitrary.  Therefore the coordinate sum is
\[
\sum_B q(R,B).
\]
Terms with \(B\not\subseteq P_B\) are zero by
\(\noderef{FixedSetProductModelMassFn}\), while for \(B\subseteq P_B\) the
definition of \(q\) factors as
\[
p_0^{|R|}(1-p_0)^{|P_R|-|R|}
\cdot p_0^{|B|}(1-p_0)^{|P_B|-|B|}.
\]
The second factor sums to \(1\) by the normalization component of
\(\noderef{FixedSetProductModelMass}\).  Thus the red outside coordinate
contributes
\[
p_0^{|D_R(x)|}(1-p_0)^{|P_R|-|D_R(x)|}.
\]
For a red coordinate \(x\in P_R\), the map \(D(v)\) imposes no condition on
the coordinate pair, so the whole coordinate sum is
\[
\sum_{(R,B)}q(R,B)=1
\]
again by \(\noderef{FixedSetProductModelMass}\).

The blue coordinates are symmetric.  For \(y\notin P_B\), the second component
is fixed to \(B=D_B(y)\), the first component is summed out, and the
coordinate contribution is
\[
p_0^{|D_B(y)|}(1-p_0)^{|P_B|-|D_B(y)|}.
\]
For \(y\in P_B\), the entire coordinate pair is unrestricted and contributes
\(1\).

Multiplying all coordinate contributions, the fiber mass for valid \(D\) is
exactly
\[
\nu(D)=
\prod_{x\notin P_R}p_0^{|D_R(x)|}(1-p_0)^{|P_R|-|D_R(x)|}
\prod_{y\notin P_B}p_0^{|D_B(y)|}(1-p_0)^{|P_B|-|D_B(y)|}.
\]
Substituting the invalid-data zero calculation and the valid-data fiber
calculation into the fiber decomposition gives
\[
\sum_{\{v:H(D(v))\}}\prod_iq(v_i)
=
\sum_{\{D\in\mathcal D:H(D)\}}\nu(D),
\]
as required.
\end{proof}
